\begin{textsample}{NEG Dim 7 – human – Score -2.00 – mg/mg\_ef\_linguagens\_ed\_fis\_3\_p2.txt}  \label{ex:f7_neg_020}
Isso significa, portanto, que devemos valorizar as estratégias \textbf{educativas} orientadas pela experimentação, pela vivência, pela prática e pela fruição e que devemos superar entendimentos que promovem a dicotomia teoria/prática, construindo uma práxis \textbf{educativa} que valorize as especificidades das ações pedagógicas da Educação Física escolar.
Para alcançarmos esse fim, é preciso trabalhar de maneira integrada, é preciso pensar, agir e sentir sem hierarquizar e classificar essas formas de interação com o mundo, com o outro e com o conhecimento.
Assumir e investir energia na defesa das especificidades das práticas escolares da Educação Física, valorizando sua importância para a formação humana integral é também uma forma de legitimar a presença desse componente na organização dos tempos e espaços escolares.
Nesse sentido, são inegáveis as expectativas vinculadas à organização de festas, apresentações, torneios, gincanas, entre tantos outros.
Elas podem e devem ser contempladas, no entanto, em sintonia e subordinação com o planejamento e a organização das práticas cotidianas que irão possibilitar o desenvolvimento das habilidades e competências aqui apresentadas.

% matched lemmas: educativo
\end{textsample}
